\begin{table}[t]
    \centering
    \caption{Comparison of optimization and convergence guarantee of deep hashing methods.}
    \label{Table: summary of deep hashing algorithms}
    \resizebox{\textwidth}{!}{
\begin{tabular}{cccccc}
\toprule
\textbf{Type} & \textbf{Method} & \textbf{Model} 
  & \textbf{Quantization Loss} & \textbf{Optimization} & \textbf{Complexity} \\
\midrule
    \multirow{3}{*}{\centering Discrete} 
    & DDSH \cite{liDeepSupervisedHashing} & $\min L_s(\bb) + Q(\bu, \bb)$ & \centering $\|\cdot\|^2$ & BQP & \multirow{12}{*}{\centering ---} \\
    \cmidrule(lr){2-5}
    & DSDH \cite{li2017deep}& \multirow{2}{*}{\centering $\min L_s(\bu) + Q(\bu, \bb)$} & \multirow{2}{*}{\centering $\|\cdot\|^2$} & $sgn$ update &  \\ 
    & DPSH \cite{liFeatureLearningBased2016} & &  & DCC &  \\ 
    \cmidrule(lr){1-5}
    \multirow{9}{*}{\centering Continuous}
    & DCH \cite{caoDeepCauchyHashing2018} &\multirow{4}{*}{\centering $\min L_s(\bu) + h(\bu)$} & $\log \left( 1 + \frac{\text{dist}_{\text{h}}(|\bu|, \mathbf{1})}{\gamma} \right)$ & \multirow{4}{*}{\centering SGDM} &  \\
    & DNNH \cite{laiSimultaneousFeatureLearning2015} & & $\mathcal{I}_{[-1,1]^d}(\bu)$ & &  \\
    & DSH \cite{liuDeepSupervisedHashing} & & $\| | \bu| - \mathbf{1}\|_1$ &  &  \\ 
    & DTSH \cite{wang2016DTSH} & & $\| | \bu| - \mathbf{1}\|_2^2$ &  &  \\ 
    & DHN \cite{zhuDeepHashingNetwork2016} & & $\sum_i\log(\cosh(|u_i|-1))$&  &  \\ 
    \cmidrule(lr){2-5}
    & MDSHC \cite{wang2023deep} & \scriptsize $\min L_c(\bu) + L_s (\bu)+ h(\bu)$ & $\| | \bu| - \mathbf{1}\|_1$ & RMSProp &  \\ 
    \cmidrule(lr){2-5}
    & HashNet \cite{Cao2017HashNetDL} & \multirow{2}{*}{\centering $\min L_s(\bu)$} & \multirow{2}{*}{\centering -} & SGDM &  \\     
    & OrthoHash \cite{hoeOneLossAll2021} & & & Adam &  \\
    % \cmidrule(lr){2-5}
    % & SBPG \cite{wang2024bregman} & \scriptsize $\min L_s(\bu) + h(\bu)$ & Weakly convex $h$ & SBregman PG & $\mathcal{O}(\varepsilon^{-4})$ \\
    \cmidrule(lr){2-6}
   & \textbf{DualHash-StoM} & \multirow{2}{*}{\small $ \min L_s(\bu) + Q(\bu,\bb) + h(\bb)$} & \multirow{2}{*}{$\| | \bb| - \mathbf{1}\|_1$} & \multirow{2}{*}{Sto Primal-Dual} & $\mathcal{O}(\varepsilon^{-4})$ \\
    & \textbf{DualHash-StoRM} & & & & $\mathcal{O}(\varepsilon^{-3})$ \\
    \bottomrule
    \end{tabular}
    }
    % \vspace{-5pt}
    \parbox{\textwidth}{\scriptsize
    \textbf{Notes}: 
    % The \textbf{Type} column indicates whether the method uses discrete or continuous optimization for quantization.
    The \textbf{Model} column shows the core optimization objective, where \(\bb\) denotes binary codes and \(\bu\) denotes continuous outputs. The \textbf{Quantization Loss} column specifies the loss (\(Q(\cdot, \cdot)\) or \(h(\cdot)\)) for quantization error, with `-' indicating no explicit loss. 
    % The \textbf{Optimization} column describes the training strategy, with discrete methods focusing on binary code updates. 
    DCC stands for Discrete Cyclic Coordinate, while BQP denotes Binary Quadratic Programming. The \textbf{Complexity} column reports the convergence rate, with `-' indicating no such analysis.
    }
     \vspace{-4ex}
\end{table}
\subsection{Related Work}
\label{Section: related work}
{We situate our contributions by reviewing related work on optimization in deep hashing, stochastic nonconvex constrained optimization, and primal-dual methods for structured composite problems.}

\paragraph{Optimization in Deep Hashing}
Improving retrieval accuracy and reducing the discrepancy between continuous representations and binary hash codes have long been central issues in deep hashing. One line of work adopts continuous relaxation and introduces explicit quantization regularization on network outputs; in particular, W-type regularizers such as $\||\mathbf{u}|-\mathbf{1}\|_1$ have been widely used to encourage binary-like activations \cite{liuDeepSupervisedHashing,wang2023deep,zhuDeepHashingNetwork2016}. Another line refines the discretization objective through greedy coding schemes, central similarity targets, or proxy-based supervision \cite{su2018greedy,yuanCentralSimilarityQuantization2020,MPDH2024,semanticHashCenters2025}. Alternative strategies directly address binary constraints by introducing auxiliary hash codes \cite{liDeepSupervisedHashing,li2017deep,liFeatureLearningBased2016}. Notably, even recent architectures retain W-type regularization; for instance, \cite{he2025mambahash} employs $\||\mathbf{u}|-\mathbf{1}\|_2^2$ as a smooth surrogate for $\||\mathbf{u}|-\mathbf{1}\|_1$ to facilitate gradient-based optimization, indicating that the practical relevance of W-type regularization has not diminished. The remaining challenge is to develop theoretically grounded optimization methods that can directly handle nonsmooth W-type regularization under nonlinear deep network mappings, which is precisely the focus of this work.
% Recent studies have further strengthened backbone architectures to improve feature quality, extending deep hashing from conventional CNNs to hybrid CNN--Transformer frameworks \cite{he2024hybridhash}, Transformer-based hashing models \cite{chen2021transhash,wu2025fth,liu2025capmh}, and state-space architectures \cite{he2025mambahash}. 
\paragraph{Stochastic Nonconvex Constrained Optimization}
{From a constrained optimization perspective, Problem~\eqref{Eq: the general nonlinear composite problem in intro} can be reformulated as a nonlinearly constrained optimization problem of the form}
\[
  \min_{x}\; f(x) + h(x) \quad \text{s.t.}\quad c(x) = 0.
\]
\blue{Several recent works have established oracle complexity results for first-order stochastic augmented-Lagrangian and penalty methods for problems of this type.} \cite{jin2022stochastic} proposed a single-loop stochastic primal-dual method based on a linearized augmented-Lagrangian function, achieving oracle complexities of $\mathcal{O}(\varepsilon^{-6})$ from an arbitrary initial point and $\mathcal{O}(\varepsilon^{-5})$ from a feasible one for locating an $\varepsilon$-KKT point. By incorporating the recursive momentum estimator of \cite{cutkoskyMomentumBasedVarianceReduction2019}, \cite{shi2025momentum} developed a momentum-based linearized augmented-Lagrangian method, improving these bounds to $\mathcal{O}(\varepsilon^{-4})$, or the nearly optimal $\mathcal{O}(\varepsilon^{-3})$ when initialized at a nearly feasible point. To certify a stronger notion of stationarity that deterministically enforces constraint violation of at most $\varepsilon$, \cite{lu2024variancereducedfirstordermethodsdeterministically} introduced a linearized quadratic-penalty method with truncated momentum. Under a regularity condition on the constraint gradients, recursive momentum yields the nearly optimal complexity $\widetilde{\mathcal{O}}(\varepsilon^{-3})$. 
% while Polyak momentum with diminishing step sizes achieves $\widetilde{\mathcal{O}}(\varepsilon^{-4})$.
\paragraph{Primal-Dual Methods for Structured Composite Problems}
In imaging science, the Primal-Dual Hybrid Gradient (PDHG) method \cite{chambolleFirstOrderPrimalDualAlgorithm2011} establishes an ergodic $\mathcal{O}(1/k)$ convergence rate for the primal-dual gap in convex settings and has found broad application in imaging inverse problems. \cite{esserZhangChan2010} proposed a unified framework for a broad class of first-order primal-dual algorithms for convex imaging problems, also achieving $\mathcal{O}(1/k)$ ergodic convergence. The stochastic extension, SPDHG \cite{chambolleSPDHG2018}, establishes almost-sure convergence and ergodic $\mathcal{O}(1/k)$ guarantees for structured finite-sum composite objectives under arbitrary sampling. \cite{chouzenouxAdjointMismatch} further showed that primal-dual methods retain convergence in the presence of adjoint mismatch, a perturbation closely related to stochastic gradient approximations; \cite{chouzenouxStochasticMM2017} developed a stochastic majorize-minimize subspace algorithm for online penalized estimation with rigorous convergence guarantees in imaging reconstruction settings. Extending to nonconvex settings, \cite{clasonAccelerationGlobalConvergence2019} established global convergence of a first-order primal-dual method under the Kurdyka--{\L}ojasiewicz condition. For nonconvex composite finite-sum objectives, \cite{guo2023preconditioned} proposed preconditioned primal-dual gradient methods with convergence analysis, while \cite{tran2022hybrid} developed a hybrid stochastic framework achieving $\mathcal{O}(\varepsilon^{-4})$ complexity. When the problem admits a block-separable structure, ADMM provides a natural decomposition framework; nonconvex ADMM convergence has been systematically established \cite{li2015global,hong2016convergence,wangGlobalConvergenceADMM2018,boct2020proximal}, and \cite{huang2019faster} established an $\mathcal{O}(\varepsilon^{-4})$ sample complexity bound for the nonconvex stochastic case. {However, the above results either rely on convex or linearly coupled structures, or address only single-block stochastic updates; the multi-block, nonconvex, nonlinearly coupled stochastic setting arising in deep hashing remains largely unaddressed.}
